% =====================================================================
\section{Preliminaries}\label{ch:prelim}
% =====================================================================

The slides open with a one-line reminder about partial derivatives and then
repeatedly say ``these partial derivatives are continuous, hence $f$ is
differentiable''. This chapter collects precisely the background behind those
sentences, so that every later step can be justified.

\subsection{Functions of several variables and level curves}

\begin{definition}[Function of two variables]
A \emph{function of two variables} $f$ assigns to each point $(x,y)$ of a set
$D\subseteq\R^2$ (the \emph{domain}) a unique real number $f(x,y)$. Its
\emph{graph} is the surface $z=f(x,y)$ in $\R^3$. Functions of three variables
$w=f(x,y,z)$ are defined the same way on sets $D\subseteq\R^3$.
\end{definition}

\begin{definition}[Level curves and level surfaces]
For a constant $k$, the \emph{level curve} of $f(x,y)$ at height $k$ is the set
of points with $f(x,y)=k$. A drawing of several level curves is a \emph{contour
map}. For $w=f(x,y,z)$ the sets $f(x,y,z)=k$ are \emph{level surfaces}.
\end{definition}

\begin{example}\label{ex:levelcurves}
Describe the level curves of $f(x,y)=4-x^2-\tfrac14 y^2$, the function of
Slides, Example~6.1.
\end{example}
\begin{solution}
Set $f(x,y)=k$: $\;4-x^2-\tfrac14y^2=k \iff x^2+\dfrac{y^2}{4}=4-k.$
\begin{itemize}
  \item If $k<4$, this is the ellipse $\dfrac{x^2}{4-k}+\dfrac{y^2}{4(4-k)}=1$ with
        semi-axes $\sqrt{4-k}$ (along $x$) and $2\sqrt{4-k}$ (along $y$).
  \item If $k=4$, only the point $(0,0)$; if $k>4$, the level set is empty.
\end{itemize}
So the graph is a ``hill'' (an elliptic paraboloid opening downward) with summit
$(0,0,4)$; the level curves are nested ellipses that get larger as $k$
decreases (Figure~\ref{fig:contour-gradient} in Chapter~\ref{ch:geometry}).
\end{solution}

\subsection{Partial derivatives}

\begin{definition}[Partial derivatives]\label{def:partials}
Let $z=f(x,y)$. The partial derivatives of $f$ at $(x_0,y_0)$ are
\[
  f_x(x_0,y_0)=\lim_{h\to0}\frac{f(x_0+h,\,y_0)-f(x_0,y_0)}{h},\qquad
  f_y(x_0,y_0)=\lim_{h\to0}\frac{f(x_0,\,y_0+h)-f(x_0,y_0)}{h},
\]
whenever these limits exist. Other notations: $f_x=\pd{f}{x}=\pd{z}{x}$,
$f_y=\pd{f}{y}=\pd{z}{y}$.
\end{definition}

\textbf{How to compute.} To find $f_x$, treat $y$ as a constant and
differentiate with respect to $x$ using ordinary single-variable rules; for
$f_y$, treat $x$ as a constant.

\textbf{Geometric meaning.} $f_x(x_0,y_0)$ is the slope, at $x=x_0$, of the
curve cut from the surface $z=f(x,y)$ by the vertical plane $y=y_0$; it is the
rate of change of $z$ as we move in the direction of $\vi=\ip{1,0}$.
Similarly $f_y(x_0,y_0)$ is the rate of change of $z$ in the direction of
$\vj=\ip{0,1}$. The whole of the next chapter is about replacing $\vi$ and
$\vj$ by an \emph{arbitrary} unit vector $\vu$.

\begin{example}\label{ex:partials}
Let $f(x,y)=x^3+x^2y^3-2y^2$. Find $f_x(2,1)$ and $f_y(2,1)$.
\end{example}
\begin{solution}
\step{1}{Differentiate in $x$ ($y$ constant).} $f_x(x,y)=3x^2+2xy^3$.
\step{2}{Differentiate in $y$ ($x$ constant).} $f_y(x,y)=3x^2y^2-4y$.
\step{3}{Evaluate.}
$f_x(2,1)=3(4)+2(2)(1)=16$ and $f_y(2,1)=3(4)(1)-4(1)=8$.

Interpretation: at $(2,1)$ the surface rises with slope $16$ in the
$x$-direction and with slope $8$ in the $y$-direction.
\end{solution}

\begin{example}\label{ex:partials2}
Find all first partial derivatives of
(a) $f(x,y)=y\ln x+xy^2$ \quad (b) $g(x,y,z)=e^{xy}\sin z$ \quad
(c) $h(x,y)=\dfrac{x}{x^2+y^2}$.
\end{example}
\begin{solution}
\textbf{(a)} $f_x=\dfrac{y}{x}+y^2$ (since $\frac{\dd}{\dd x}\ln x=\frac1x$),
\quad $f_y=\ln x+2xy$. (These are used in Slides, Example~6.3.)

\textbf{(b)} Chain rule on $e^{xy}$:
$g_x=ye^{xy}\sin z,\quad g_y=xe^{xy}\sin z,\quad g_z=e^{xy}\cos z.$

\textbf{(c)} Quotient rule:
\[
h_x=\frac{(x^2+y^2)(1)-x(2x)}{(x^2+y^2)^2}=\frac{y^2-x^2}{(x^2+y^2)^2},\qquad
h_y=\frac{(x^2+y^2)(0)-x(2y)}{(x^2+y^2)^2}=\frac{-2xy}{(x^2+y^2)^2}.
\]
\end{solution}

\subsection{Differentiability}

For one variable, ``differentiable'' just means the derivative exists. For two
variables, the mere existence of $f_x$ and $f_y$ is \emph{not} enough for the
formulas of this course; we need the stronger notion below.

\begin{definition}[Differentiable function]
$f(x,y)$ is \emph{differentiable} at $(x_0,y_0)$ if $f_x(x_0,y_0)$ and
$f_y(x_0,y_0)$ exist and
\[
  f(x_0+\Delta x,\,y_0+\Delta y)-f(x_0,y_0)
  = f_x(x_0,y_0)\,\Delta x+f_y(x_0,y_0)\,\Delta y+\varepsilon_1\Delta x+\varepsilon_2\Delta y,
\]
where $\varepsilon_1,\varepsilon_2\to0$ as $(\Delta x,\Delta y)\to(0,0)$.
Geometrically: near $(x_0,y_0)$ the surface is well approximated by its
tangent plane $z=f(x_0,y_0)+f_x(x_0,y_0)(x-x_0)+f_y(x_0,y_0)(y-y_0)$.
\end{definition}

Checking this definition directly is tedious. In practice we use the following
sufficient condition, which is exactly what the slides invoke.

\begin{theorem}[Continuous partials imply differentiability]\label{thm:C1}
If $f_x$ and $f_y$ exist near $(x_0,y_0)$ and are continuous at $(x_0,y_0)$,
then $f$ is differentiable at $(x_0,y_0)$. The same holds for functions of
three variables.
\end{theorem}

\begin{remark}[How to use it on an exam]
Polynomials, exponentials, $\sin$, $\cos$, and their sums, products, quotients
(away from zeros of the denominator) and compositions have continuous partial
derivatives on their domains. So a one-line justification such as
``$f_x=2x\sin2y$ and $f_y=2x^2\cos2y$ are continuous on $\R^2$, hence $f$ is
differentiable'' is enough, and you should write it, because the theorem of
Chapter~\ref{ch:dirder} \emph{requires} differentiability.
\end{remark}

\begin{warning}
Existence of $f_x$ and $f_y$ alone does \emph{not} imply differentiability.
Problem~\ref{prob:nondiff} in Chapter~\ref{ch:practice} gives a function with
$f_x(0,0)=f_y(0,0)=0$ whose directional derivatives at the origin do not obey
the formula of Theorem~\ref{thm:Du-formula}.
\end{warning}

\subsection{The chain rule along a curve}

\begin{theorem}[Chain rule, Case 1]\label{thm:chain}
If $z=f(x,y)$ is differentiable and $x=x(t)$, $y=y(t)$ are differentiable
functions of $t$, then $z$ is a differentiable function of $t$ and
\[
  \frac{\dd z}{\dd t}=\pd{f}{x}\frac{\dd x}{\dd t}+\pd{f}{y}\frac{\dd y}{\dd t}.
\]
For $w=f(x,y,z)$ along $x(t),y(t),z(t)$ there is one more term:
$\frac{\dd w}{\dd t}=f_x x'(t)+f_y y'(t)+f_z z'(t)$.
\end{theorem}

\begin{example}\label{ex:chain}
Let $z=x^2y+3xy^4$ with $x=\sin2t$, $y=\cos t$. Find $\dfrac{\dd z}{\dd t}$ at $t=0$.
\end{example}
\begin{solution}
\step{1}{Partials.} $z_x=2xy+3y^4$, $z_y=x^2+12xy^3$.
\step{2}{Point.} At $t=0$: $x=\sin0=0$, $y=\cos0=1$, so $z_x=0+3=3$ and $z_y=0+0=0$.
\step{3}{Derivatives of $x,y$.} $x'(t)=2\cos2t$, $y'(t)=-\sin t$, so $x'(0)=2$, $y'(0)=0$.
\step{4}{Chain rule.} $\dfrac{\dd z}{\dd t}\Big|_{t=0}=3\cdot2+0\cdot0=6.$
\end{solution}

The chain rule is the engine behind both the formula for directional
derivatives (Theorem~\ref{thm:Du-formula}) and the fact that the gradient is
perpendicular to level curves (Theorem~\ref{thm:grad-perp}).

\subsection{Unit vectors, angles and the dot product}

\begin{itemize}
  \item The \emph{length} of $\vv=\ip{v_1,v_2}$ is $\norm{\vv}=\sqrt{v_1^2+v_2^2}$
        (in $\R^3$: $\sqrt{v_1^2+v_2^2+v_3^2}$).
  \item A \emph{unit vector} has length $1$. Any nonzero $\vv$ can be
        \emph{normalized}: $\vu=\dfrac{\vv}{\norm{\vv}}$ is the unit vector with the
        same direction as $\vv$.
  \item The unit vector making angle $\theta$ with the positive $x$-axis is
        $\vu=\ip{\cos\theta,\sin\theta}$.
  \item The vector from $P(x_1,y_1)$ to $Q(x_2,y_2)$ is
        $\vecPQ=\ip{x_2-x_1,\,y_2-y_1}$.
  \item Dot product: $\vb a\cdot\vb b=a_1b_1+a_2b_2(+a_3b_3)
        =\norm{\vb a}\norm{\vb b}\cos\theta$, where $\theta\in[0,\pi]$ is the angle
        between $\vb a$ and $\vb b$. In particular $\vb a\cdot\vb b=0$ iff
        $\vb a\perp\vb b$.
\end{itemize}

\begin{example}\label{ex:unit}
(a) Normalize $\vv=3\vi-4\vj$. (b) Find the unit vector in the direction from
$P(-\tfrac34,0)$ to $Q(0,1)$. (c) Write the unit vector at angle
$\theta=\tfrac{2\pi}{3}$.
\end{example}
\begin{solution}
\textbf{(a)} $\norm{\vv}=\sqrt{9+16}=5$, so $\vu=\ip{\tfrac35,-\tfrac45}$.

\textbf{(b)} $\vecPQ=\ip{0-(-\tfrac34),\,1-0}=\ip{\tfrac34,1}$ and
$\norm{\vecPQ}=\sqrt{\tfrac{9}{16}+1}=\sqrt{\tfrac{25}{16}}=\tfrac54$, so
$\vu=\ip{\tfrac34,1}\big/\tfrac54=\ip{\tfrac35,\tfrac45}$.

\textbf{(c)} $\vu=\ip{\cos\tfrac{2\pi}{3},\sin\tfrac{2\pi}{3}}=\ip{-\tfrac12,\tfrac{\sqrt3}{2}}$.
\end{solution}

\begin{remark}
Clearing fractions before normalizing saves work: $\ip{\tfrac34,1}$ has the same
direction as $4\ip{\tfrac34,1}=\ip{3,4}$, whose length is $5$. Hence
$\vu=\ip{\tfrac35,\tfrac45}$ immediately.
\end{remark}
